\subsection{Cost accounting, historical placement, and formal certification}\label{sec:realize:cost}

\begin{definition}[Affine compilation parameters]\label{def:realize:complexity}
Let $s$ be the number of distinct secret field
elements registered in one stage, and let
$E$ be the number of field-affine equations
verified after registration.
For the $e$-th equation and $k$-th
output coordinate, write $m_{e,k}$
for the number of nonzero binary
coefficients in that parity row.
Write $B:=sn$ for the number
of committed binary coordinates.
\end{definition}

\begin{lemma}[Exact bit commitment and parity proof counts]\label{lem:realize:count}
The compiler registers $B=sn$ Pedersen
commitments and, if every coordinate is
registered once, executes $B$ binary
bitness OR proofs.
It executes exactly $nE$ parity rows,
except that a zero-length row may be
checked directly and omitted.
Each nonempty row's OR proof has
\[
 d_{e,k}=|\mathcal S_{m_{e,k},b_{e,k}}|
 \leq\left\lceil\frac{m_{e,k}+1}{2}\right\rceil
\]
candidate branches.
\end{lemma}
\begin{proof}
Every field element has exactly
$n$ binary basis coordinates,
so $s$ elements supply $sn=B$
commitments and one bitness
proof per coordinate.
Theorem~\ref{thm:realize:affine-to-parity}
produces one parity equation for
each of the $n$ output coordinates
of each of the $E$ field equations.
For a row with $m$ inputs,
the allowed integers in $[0,m]$
of one specified parity form
one of the two parity classes.
The larger class has
$\lceil(m+1)/2\rceil$ elements,
giving the branch bound.
The empty-row handling follows
Corollary~\ref{cor:realize:empty-row}.
\end{proof}

\begin{lemma}[Linear branch-size upper bound]\label{lem:realize:branch-upper}
The total number of candidate
branches across all nonempty
parity OR proofs is at most
\[
 \boxed{
 \sum_{e=1}^{E}\sum_{k=0}^{n-1}
 \left\lceil\frac{m_{e,k}+1}{2}\right\rceil
 \leq nE\left\lceil\frac{sn+1}{2}\right\rceil.}
\]
\end{lemma}
\begin{proof}
The first inequality follows
from Lemma~\ref{lem:realize:count}
applied rowwise, and zero rows
contribute no actual proof branches.
Every row weight is at most
$sn$ by Lemma~\ref{lem:realize:row-weight}.
The ceiling function is monotone,
so each term in the double sum
is at most $\lceil(sn+1)/2\rceil$.
There are $nE$ such terms,
giving the second inequality.
\end{proof}

\begin{lemma}[Two group elements per branch in a raw OR proof]\label{lem:realize:or-communication}
An uncompressed $d$-branch
Schnorr OR proof transmits $d$
group elements in its first
message and $2d$ elements of
$\mathbb Z_P$ in its responses,
in addition to one verifier
challenge in $\mathbb Z_P$.
\end{lemma}
\begin{proof}
Definition~\ref{def:realize:or}
sends the vector $(T_j)_{j=1}^d$
of $d$ group elements, followed
by the vector $(e_j,z_j)$ of
two scalars per branch.
The verifier supplies one
global challenge $e$.
No compression or derived
challenge optimization is
assumed.
\end{proof}

\begin{proposition}[Raw compiler overhead]\label{prop:realize:not-succinct}
The raw compiler commits to every
secret field coordinate separately.
Its parity proofs also grow with
the number of binary row entries.
Consequently the construction
in this chapter \emph{does not}
establish a polylogarithmic-size
zero-knowledge PCS, even if
the original KBFold commitment
scheme has succinct opening proofs.
\end{proposition}
\begin{proof}
Lemma~\ref{lem:realize:count}
requires $sn$ separate bit
commitments and bitness proofs
in the direct implementation.
Lemma~\ref{lem:realize:branch-upper}
counts a candidate set for
each parity row, and
Lemma~\ref{lem:realize:or-communication}
requires a group message
and scalar responses for
every candidate branch.
These costs are proportional
to explicit field-coordinate
and row representations;
nothing in the construction
compresses them to a
polylogarithmic dependence
on the witness length.
\end{proof}

\begin{definition}[Chapter 10 Lean module contract]\label{def:realize:lean-map}
The formalization splits into:
\begin{enumerate}[label=\textup{(\roman*)}]
\item exact finite-field affine-to-binary
coordinate matrices;
\item prime-order cyclic group algebra
and Pedersen hiding/binding;
\item Schnorr and OR-protocol
message equations, HVZK simulators
and two-transcript extraction;
\item integer-lift parity candidate
sets and wrong-parity
discrete-log reduction;
\item compilation of all Chapter~9
affine equations and packing links;
\item finite adaptive public-core plus
honest-sigma proof-transcript simulation;
\item a parameterized early-extraction
contract and its conditional
outer soundness implication.
\end{enumerate}
\end{definition}

\begin{remark}[Prior work and originality]
The commitment construction and
its hiding/binding theorem come
from Pedersen~\cite{Pedersen1991}.
The logarithm proof is Schnorr's
\cite{Schnorr1991}, while OR
proofs of partial knowledge
are classical
\cite{CramerDamgardSchoenmakers1994}.
The affine-to-parity application
here is an explicit composition
in our Boolean-kernel
Ring Switching notation,
not a claim to have invented
bit commitments or OR proofs.
The target PCS and its
post-quantum motivation are
from the original KBFold
manuscript~\cite{MkhidaKBFold2026}.
\end{remark}

\begin{remark}[Main remaining theorem]
The next goal is a malicious-verifier
simulator and a sound extractor for
the bound commitments. Both must
remain secure under adaptive outer
challenges.
After that, one must bridge
these characteristic-two
bit commitments to KBFold's
odd-characteristic, authenticated
folding-oracle PCS.
\end{remark}
